% Header aus der Vorlage
\documentclass[a4paper,11pt, footheight=26pt]{scrartcl}
\usepackage[head=23pt]{geometry}	% head=23pt umgeht Fehlerwarnung, dafür größeres "top" in geometry
\geometry{a4paper, top=30mm, left=20mm, right=20mm, bottom=22mm,headsep=10mm, footskip=12mm}

% Input inkl. Umlaute, Silbentrennung
\usepackage[T1]{fontenc}
\usepackage[utf8]{inputenc}
\usepackage[ngerman]{babel}
\usepackage{csquotes}	% Anführungszeichen

% HTW Corporate Design: Arial (Helvetica)
\usepackage{helvet}
\renewcommand{\familydefault}{\sfdefault}

% Style-Aufhübschung
\usepackage{soul, color}	% Kapitälchen, Unterstrichen, Durchgestrichen usw. im Text
\usepackage{scrlayer-scrpage}	% Kopf-/Fußzeile
%\usepackage{titleref}
\usepackage[perpage]{footmisc}	% Fußnotenzählung Seitenweit, nicht Dokumentenweit

% Mathe usw.
\usepackage{amssymb}
\usepackage[fleqn]{amsmath}	% fleqn: align-Umgebung rechtsbündig
\usepackage{xcolor}
\usepackage{esint}	% Schönere Integrale, \oiint vorhanden
\everymath=\expandafter{\the\everymath\displaystyle}	% Mathe Inhalte werden weniger verkleinert
\usepackage{wasysym}	% mehr Symbole, bspw \lightning

% tikz usw.
\usepackage{tikz}
\usepackage{pgfplots}
\pgfplotsset{compat=1.11}	% Umgeht Fehlermeldung
\usetikzlibrary{graphs}
%\usetikzlibrary{through}	% ???
\usetikzlibrary{arrows}
\usetikzlibrary{arrows.meta}	% Pfeile verändern / vergrößern: \draw[-{>[scale=1.5]}] (-3,5) -> (-3,3);
\usetikzlibrary{automata,positioning} % Zeilenumbruch im Node node[align=center] {Text\\nächste Zeile}
\usetikzlibrary{patterns}	% Schraffierte Füllung
\tikzstyle{reverseclip}=[insert path={	% Inverser Clip \clip
	(current page.north east) --
	(current page.south east) --
	(current page.south west) --
	(current page.north west) --
	(current page.north east)}
% Nutzen: 
%\begin{tikzpicture}[remember picture]
%\begin{scope}
%\begin{pgfinterruptboundingbox}
%\draw [clip] DIE FLÄCHE, IN DER OBJEKT NICHT ERSCHEINEN SOLL [reverseclip];
%\end{pgfinterruptboundingbox}
%\draw DAS OBJEKT;
%\end{scope}
%\end{tikzpicture}
]	% Achtung: dafür muss doppelt kompliert werden!
\usepackage{graphpap}	% Grid für Graphen

% Tabular
\usepackage{longtable}	% Große Tabellen über mehrere Seiten
\usepackage{multirow}	% Multirow/-column: \multirow{2[Anzahl der Zeilen]}{*[Format]}{Test[Inhalt]} oder \multicolumn{7[Anzahl der Reihen]}{|c|[Format]}{Test2[Inhalt]}
\renewcommand{\arraystretch}{1.3} % Tabellenlinien nicht zu dicht
\usepackage{colortbl}
\arrayrulecolor{gray}	% heller Tabellenlinien

% Nützliches
\usepackage{verbatim}	% u.a. zum auskommentieren via \begin{comment} \end{comment}
\usepackage{tabto}	% Tabs: /tab zum nächsten Tab oder /tabto{.5 \CurrentLineWidth} zur Stelle in der Linie
\NumTabs{6}	% Anzahl von Tabs pro Zeile zum springen
\usepackage{listings} % Source-Code mit Tabs
\usepackage{lstautogobble} 
\usepackage{enumitem}	% Anpassung der enumerates
\setlist[enumerate,1]{label=\arabic*.)}	% global andere Enum-Items
\usepackage{letltxmacro} % neue Definiton von Grundbefehlen
% Nutzen:
%\LetLtxMacro{\oldemph}{\emph}
%\renewcommand{\emph}[1]{\oldemph{#1}}

% Einrichtung von lst
\lstset{
basicstyle=\ttfamily, 
mathescape=true, 
escapeinside=||, 
autogobble, 
tabsize=2,
basicstyle=\footnotesize\sffamily\color{black},
frame=single,
rulecolor=\color{lightgray},
numbers=left,
numbersep=5pt,
%basicstyle=\footnotesize\sffamily\color{black},
commentstyle=\color{gray},
%frame=single,
%numbers=left,
%numbersep=5pt,
numberstyle=\tiny\color{gray},
keywordstyle=\color{green},
%showspaces=false,
showstringspaces=false,
stringstyle=\color{orange},
tabsize=2,
literate=%
    {Ö}{{\"O}}1
    {Ä}{{\"A}}1
    {Ü}{{\"U}}1
    {ß}{{\ss}}1
    {ü}{{\"u}}1
    {ä}{{\"a}}1
    {ö}{{\"o}}1
    {~}{{\textasciitilde}}1
}

% BibTeX
\usepackage[backend=bibtex, bibencoding=ascii]{biblatex}	% BibTeX
\usepackage{makeidx}
%\makeglossary
%\makeindex

% Grafiken
\usepackage{graphicx}
\usepackage{epstopdf}	% eps-Vektorgrafiken einfügen

% pdf-Setup
\usepackage{pdfpages}
\usepackage[bookmarks,%
bookmarksopen=false,% Klappt die Bookmarks in Acrobat aus
colorlinks=true,%
linkcolor=black,%
citecolor=red,%
urlcolor=green,%
]{hyperref}

% Titel, Autor usw. werden vor dem Anfang des Dokuments in einem Rutsch definiert…
\newcommand{\DTitel}[1]{\newcommand{\Dokumententitel}{#1}}
\newcommand{\DUntertitel}[1]{\newcommand{\Dokumentenuntertitel}{#1}}
\newcommand{\DAutor}[1]{\newcommand{\Dokumentenautor}{#1}}
\newcommand{\DNotiz}[1]{\newcommand{\Dokumentennotiz}{#1}}
% … Deswegen folgendes erst Nach Dokumentenbeginn ausführen:
\AtBeginDocument{
	\hypersetup{
		pdfauthor={\Dokumentenautor},
		pdftitle={HTW Dresden | \Dokumententitel - \Dokumentenuntertitel},
	}
	\automark[section]{section}
	\automark*[subsection]{subsection}
	\pagestyle{scrheadings}
	\ihead{\includegraphics[height=1.7em]{../LaTeX_master/HTW-Logo.eps}}
	\ohead{\Dokumententitel}
	\ofoot{\footnotesize{\textcolor{darkgray}{Mitschrift von\\ \Dokumentenautor}}}
	% Titelseite
	\title{\includegraphics[width=0.35\textwidth]{../LaTeX_master/HTW-Logo.eps}\\\vspace{0.5em}
	\Huge\textbf{\Dokumententitel} \\\vspace*{0,5cm}
	\Large \Dokumentenuntertitel \\\vspace*{4cm}}
	\author{\textcolor{darkgray}{Mitschrift von \Dokumentenautor} \vspace*{1cm}\\\Dokumentennotiz}
}

%% EIGENE BEFEHLE

%Farbdefinitionen
\definecolor{red}{RGB}{180,0,0}
\definecolor{green}{RGB}{75,160,0}
\definecolor{blue}{RGB}{0,75,200}
\definecolor{orange}{RGB}{255,128,0}
\definecolor{yellow}{RGB}{255,245,0}
\definecolor{purple}{RGB}{75,0,160}
\definecolor{cyan}{RGB}{0,160,160}
\definecolor{brown}{RGB}{120,60,10}

% Textfarbe ändern
\newcommand{\tred}[1]{\textcolor{red}{#1}}
\newcommand{\tgreen}[1]{\textcolor{green}{#1}}
\newcommand{\tblue}[1]{\textcolor{blue}{#1}}
\newcommand{\torange}[1]{\textcolor{orange}{#1}}
\newcommand{\tyellow}[1]{\textcolor{yellow}{#1}}
\newcommand{\tpurple}[1]{\textcolor{purple}{#1}}
\newcommand{\tcyan}[1]{\textcolor{cyan}{#1}}
\newcommand{\tbrown}[1]{\textcolor{brown}{#1}}

% Umstellen der Tabellen Definition
\newcommand{\mpb}[1][.3]{\begin{minipage}{#1\textwidth}\vspace*{3pt}}
\newcommand{\mpe}{\vspace*{3pt}\end{minipage}}

\newcommand{\resultul}[1]{\underline{\underline{#1}}}
\newcommand{\parskp}{$ $\\}	% new line after paragraph
\newcommand{\corr}{\;\widehat{=}\;}
\newcommand{\mdeg}{^{\circ}}

\newcommand{\nok}[2]{\begin{pmatrix}#1\\#2\end{pmatrix}}	% n über k

\bibliography{../Literatur/HTW_Literatur.bib}

% Definition von Titel, Autor usw.
\DTitel{Programmierung II}
\DUntertitel{Vorlesungsskript}
\DAutor{Falk-Jonatan Strube}
\DNotiz{Vorlesung von Dr. Arnold Beck}

\begin{document}

\maketitle
\newpage
\tableofcontents
\newpage

\section*{Hinweise}
Unterlagen unter:
\begin{lstlisting}[language=bash]
cd /home/rex/fi1/nestler/Programmierung_II_2016/
\end{lstlisting}
\paragraph{Compiler} 
\begin{itemize}
\item Intel i16, i13 (für Linux oder Visual Studio) \url{www.hocomputer.de} (kostenpflichtig)
\item gcc 5.3, 4.85 \url{gcc.gnu.org}
\end{itemize}
Zugriff auf Windows-Programme (Visual Studio 2013) in Linux-Laboren:
\begin{lstlisting}[language=bash]
rdesktop -f its56 # oder its59
\end{lstlisting}
Empfohlene Literatur: Breymann\cite{breymann2009c++}

\chapter{C++}
\section{Ein- und Ausgabe}
(siehe Folie CPP\_01\_stdio)
\begin{lstlisting}[language=C++]
#include <iostream>	// alternativ zu <stdio.h> in C
using namespace std;	// namespace: für bestimmte Abkürzungen (bpsw. cout anstatt std::cout)
// Hinweis: "::" zeigt, dass das davorstehende "static" ist (hier: std)

class integer {}; // class: Vergleichbar mit typedef

int main() {
	integer i0;	// i0: Instanz bzw. Objekt der Klasse integer
	cin.get();	// Eingabe abwarten
}
\end{lstlisting}
(vgl. integer.cpp)
\begin{lstlisting}[language=C++]
#include <iostream>
using namespace std;

class integer { // int - Variable in class verpacken
//	private:    // private ist default
		int i;  // this->i bzw. (*this).i
		// private nur für andere Klassen, andere Instanzen der gleichen Klasse können drauf zugreifen
	public: // wenn nichts steht, wird automatisch das vorherige angenommen. Hier: private
		integer(int i=0):i(i){  // Konstruktor und Defaultkonstruktor
			// i=0 default Wert, wenn keiner Angegeben
			// :i(i) übergebenes i wird dem i der Klasse zugewiesen: i_1(i_2) i_1 ist this->i, und i_2 ist das übergebene i
			cout<<"integer-Objekt i = "<<this->i<<endl; 
		}

		int get(){ return i; }

		void set(int i=0){ this->i = i; }

		// statische Methode: aufrufbar ohne Instanziierung der Klasse
		static integer add(integer i1, integer i2){  // Wertkopien von i1 und i2
			// return integer(i1.i + i2.i); // alternativ und explizit: Konstruktor-Aufruf
			// return erstellt eine Kopie (mit Konstruktur erstellt)!
			// return i1.get()+i2.get();    // Umwandlung int nach integer, Aufruf Konstruktor implizit
			return i1.i + i2.i;
			// i1.i Möglich, da innerhalb der Klasse integer und somit privates i sichtbar
		}
};

auto max(int x, int y) -> int { return x>y ? x : y; }	// Lambda-Funktion
// auto: Rückgabetyp ergibt sich aus dem Kontext bzw. über das "-> int"

template<typename Typ1, typename Typ2>	// Weiterentwicklung Makro: wählt automatisch Typ aus
auto quotient(Typ1 a, Typ2 b) -> decltype(a/b) { return a/b; }

auto main() -> int {
	auto k = 0;        // C++11: da 0 vom Typ int ist auch k vom Typ int
	decltype(k) j = 5; // C++11: da k vom Typ int ist auch j vom Typ int
	char *c = nullptr; //C++ 11: Zeigerliteral
	int *ip = NULL;

	integer i0(5), i1=4;     // 2 (alternative) Initialisiierungen von Objekten
	// i1=4 nur möglich, wenn 1 Parameter gefordert ist.
	cout<<"i0.i = "<<i0.get()<<endl;
	// cout im Vergleich zu printf() typsicher.
	cout<<"i0.i + i0.i = "<<integer::add(i0, i0).get()<<endl; // Aufruf static-Methode add
	integer i3 = integer::add(i0, i0);                        // Initialisierung von i3
	cout<<"i3.i     = "<<i3.get()<<endl;
	i0.set(22);
	cout<<"i0.i     = "<<i0.get()<<endl;
	cout<<"max(3,5) = "<<max(3,5)<<endl;
	cout<<"5/3      = "<<quotient(5, 3)<<endl;
	cout<<"5.0/3.0  = "<<quotient(5.0, 3.0)<<endl;
	cout<<"b / 1    = "<<quotient('b', '1')<<endl;
	cin.get();
}
\end{lstlisting}

(vgl iostream.pdf)
\begin{itemize}
\item nach jeder cin Eingabe: „cin.clear();“, damit Fehler ignoriert werden um weiter cin's abhandeln zu können (vgl. robust\_ea)
\end{itemize}
Einlesen:
\begin{lstlisting}[language=C++]
char sc;
cout << "sc=";
cin >> sc;
cin.clear();	// clear, um die Eingabezeile freizumachen, damit man nicht an Falscher Eingabe hängen bleibt
cin.ignore(INT_MAX, '\n');	// braucht #include <limits.h>
cout << "sc" << dec << (int) sc << endl;

// alternativ:
char vb[128];
cout << "s=";
cin.getline(vb, sizeof(vb), '\n');	// lesen als String, dann wieder umwandeln (liest auch Leerzeichen ein)
sc = atoi(vb);	// braucht #include <cstdlib>

// alternative zu getline:
cin.get(...);	// lässt aber \n im Strom
cin.get();

// alternativ
cin >> setw(sizeof(vb)) >> vb;	// verhindert Überlauf
sc = atoi(vb);

// alternativ
String s;	// braucht #include <string>
size_t ie=0;
cin >> s;
unsigned int ni = stoi(s, &ie, 10);

// alternativ
getline(cin, s, '\n');
double d = stod(s, &ie);

// zum compilieren: g++ p2a1.cpp -std=c++11 -o a.out
\end{lstlisting}

robust\_ea1.cpp:
\begin{lstlisting}[language=C++]
	do {
		cout<<"d = "; cin>>d; 	// einlesen
		if(cin.eof()) break;		// break bei Strg+D oä.
		if(cin.fail() || (cin.peek() != '\n')){		// ist nächstes Zeichen ungültig?
			cin.clear(); cin.ignore(INT_MAX,'\n');	// Strom zurücksetzen und zum \n gehen
			continue;
		}
		break; % Schleife verlassen, wenn korrekte Eingabe
	} while(true);
	
	if(cin.eof()){ cin.clear(); cout<<"eof\n"; }
	else {
		cin.clear(); cin.ignore(INT_MAX,'\n');
		cout<<"Wert d = "<<d<<endl;
	}
	cin.ignore();
\end{lstlisting}

\section{Defaultargumente}
Defaultargumente müssen immer von rechts angefangen definiert sein:
\begin{lstlisting}[language=C++]
myFunc(int i = 5, int j = 7)	// korrekt
myFunc(int i, int j = 7)			// korrekt
myFunc(int i = 5, int j)			// falsch!!!
\end{lstlisting}

\section{Überladen}
overload.pdf

Hinweis: cast auf zwei Möglichkeiten:
\begin{lstlisting}[language=C++]
int i = 5;
double d;
d = (long) i;
d = long(i);
\end{lstlisting}

\section{Typisierte Konstanten}

\section{Referenzen}
referenzen.pdf\\
Ein Speicherplatz wird mit mehreren Variablen-Namen beschrieben.

\newpage
\printbibliography
\end{document}